\documentclass[10pt,a4paper]{amsart}
\usepackage[margin=19mm]{geometry}
\usepackage{amsmath,amssymb,mathtools}
\usepackage[scr=rsfs,cal=euler]{mathalfa}
\usepackage{xfrac}
\usepackage[unicode,hidelinks,bookmarks=false]{hyperref}
\newcommand{\ie}{i.e.\ }
\newcommand{\cf}{cf.\ }
\newcommand{\resp}{respectively\ }
\newcommand{\Subsection}{\subsection}
\providecommand{\nobreakdash}{\nobreak\hbox{-}}
\setlength{\emergencystretch}{2em}
\multlinegap0pt
\usepackage{bm}
\usepackage{bbm}
\usepackage{dsfont}
\usepackage{xspace}
\usepackage{cancel}
\usepackage{mathtools}
\usepackage{amsthm}
\makeatletter
\@ifundefined{theorem}{\newtheorem{theorem}{Theorem}}{}
\@ifundefined{claim}{\newtheorem{claim}{Claim}}{}
\@ifundefined{lemma}{\newtheorem{lemma}[theorem]{Lemma}}{}
\makeatother
\newcommand{\F}{\mathbb{F}}
\newcommand{\PG}{\operatorname{PG}}
\title[On the generalized Tur\'{a}n number of the complete bipartit]{On the generalized Tur\'{a}n number of the complete bipartite graph $K\_{3,b+1}$}
\author{Jing Wang, Zixuan Yang, Junpeng Zhou}
\date{Source: arXiv:2607.01680v1 (CC BY 4.0)}
\begin{document}
\maketitle

\noindent\emph{Source and licence.} On the generalized Tur\'{a}n number of the complete bipartite graph $K_{3,b+1}$. Version arXiv:2607.01680v1 (CC BY 4.0), distributed under the Creative Commons Attribution 4.0 International licence, as recorded in the arXiv licence metadata for this exact version and shown on the abstract page https://arxiv.org/abs/2607.01680v1. Full rights notice: licenses/arxiv-2607.01680-CC-BY-4.0.txt.\par\smallskip

\noindent\textit{Editorial scope.} Counted unit: the Lemma whose source label is \texttt{lem:good-plane-count} in the article (source0.tex lines 576--578, source label \texttt{lem:good-plane-count}), delivered together with the article's own complete original proof of it (source0.tex lines 580--766, one \texttt{proof} environment of 187 lines). No mathematics is written, repaired or re-derived: statement and proof are the article's own text line for line. Required context retained uncounted: lem:double-shift (lines 238--244); lem:no-three-collinear (lines 349--351). Every definition, display and label of those lines is retained unchanged. Omitted from this package and not redistributed: 1--237, 245--348, 352--575, 579--579, 767--951. Nothing inside a retained excerpt is omitted or altered: every retained body is delivered exactly as the article's lines are written, so no omission receipt of any kind accompanies this package. Citations: the retained excerpts carry no citation command at all, so no bibliography entry of the article is transcribed and no reference is redistributed. Reference adaptation: none was needed and none was made; every reference inside the delivered unit resolves inside it, so no unresolved reference or citation is produced. Printed slips kept literal: no printed word, symbol, hypothesis or number of the counted statement or of its proof was altered. Layout and class: the article's own class is \texttt{article} with packages including amsmath, amssymb, amsthm, amsfonts, graphicx, hyperref, enumerate, mathtools; the wrapper is the lane's pinned \texttt{amsart} editorial wrapper at 10pt on A4, which restates the theorem-like environments the retained text needs and adds the article's own macro definitions that the retained text needs (\texttt{\textbackslash F}, \texttt{\textbackslash PG}), with extra wrapper packages (bm, bbm, dsfont, xspace, cancel, mathtools). The delivered numbering is the unit's own. Assets: no retained span contains a figure environment, a graphics-inclusion command, a PDF-page inclusion, a table or a TikZ picture; no image file is copied into this package, the package declares no asset dependency and no image file is redistributed with it. Licence: version arXiv:2607.01680v1 of this article is distributed under the Creative Commons Attribution 4.0 International licence, as recorded in the arXiv licence metadata for this exact version and shown on the abstract page https://arxiv.org/abs/2607.01680v1. A full rights notice is delivered at licenses/arxiv-2607.01680-CC-BY-4.0.txt. Characters: every retained excerpt is pure ASCII, so the delivered engine, XeLaTeX with its default Latin Modern text font, reports no missing character.
\begin{lemma} \label{lem:double-shift}
Let $b \geq 5$ be an integer. There exist constants $q_0=q_0(b)$ and $\eta_b > 0$ such 
that for $q \geq q_0$ and polynomial 
$\delta \in \F_q[t]$ with $\deg \delta < b$, there are at least $\eta_b q^b$ polynomials 
$f \in \mathcal{M}_b$ such that both $f$ and $f + \delta$ split into $b$ 
distinct linear factors over $\F_q$.
\end{lemma}

\begin{lemma} \label{lem:no-three-collinear}
For every $f \in \mathcal{M}_b$, the set $S_f$ contains no three distinct collinear points.
\end{lemma}

\begin{lemma} \label{lem:good-plane-count}
Fix integers $b \geq 5$ and $4 \leq a \leq b$. For all sufficiently large odd prime powers $q$, there exists a $f \in \mathcal{M}_b$ such that  the bipartite graph $G_f$ has $2q^3$ vertices and $\Omega_{a,b}(q^9)$ copies of $K_{a,b}$.
\end{lemma}

\begin{proof}
Let $\PG(5,q)$ be a 5-dimensional projective space  over $\mathbb{F}_q$  with the nondegenerate symmetric bilinear form $\beta$ that induces an orthogonal polarity. Recall the point set
\[
S_f = \left\{ s_f(x, y, z) : x, y, z \in \mathbb{F}_q \right\},
\]
where 
\[
s_f(x, y, z) = \left[ 1 : x : y : z : x^2 : y^2 - \alpha z^2 + f(x) \right],
\]
with $\alpha \in \mathbb{F}_q$ a fixed non-square and $f \in \mathcal{M}_b$. By the definition of  $G_f$, $G_f$ has exactly $2|S_f| = 2q^3$ vertices.

We first introduce a 9-parameter family of projective planes. 
For any triple of polynomials of degree at most $2$:
\[
\begin{aligned}
A(t) &= a_0 + a_1 t + a_2 t^2, \\
B(t) &= b_0 + b_1 t + b_2 t^2, \\
C(t) &= c_0 + c_1 t + c_2 t^2,
\end{aligned}
\]
we define the projective plane $\Pi(A,B,C)$ as the set of points satisfying the linear system
\[
\begin{aligned}\label{exKabK3b+1-equation7}
y &= a_0 w + a_1 x + a_2 u, \\
z &= b_0 w + b_1 x + b_2 u, \\
v &= c_0 w + c_1 x + c_2 u.
\end{aligned}
\]
Each equation is a homogeneous linear relation, and the three equations are linearly independent (each isolates a distinct coordinate $y,z,v$ with leading coefficient $1$). The solution space is therefore a $3$-dimensional vector subspace, i.e., a projective plane in $\PG(5,q)$.
Distinct triples $(A,B,C)$ yield distinct planes. Indeed, if two triples produce the same plane, the linear expressions for $y,z,v$ in terms of $w,x,u$ must coincide termwise, so all coefficients $a_i,b_i,c_i$ must be equal. Thus, the full family therefore contains $q^3 \cdot q^3 \cdot q^3 = q^9$ distinct planes.



We now establish a bijection between $S_f \cap \Pi(A,B,C)$ and the roots of a polynomial over $\mathbb{F}_q$. A point $s_f(x,y,z) \in S_f$ belongs to $\Pi(A,B,C)$ if and only if its coordinates satisfy the three plane equations. Substituting $w=1$, $u=x^2$, $v = y^2 - \alpha z^2 + f(x)$ into (\ref{exKabK3b+1-equation7}), we get
\[
\begin{aligned}
~~~~~~~~~~~~~~~~~&y = a_0 + a_1 x + a_2 x^2 = A(x), \\
&z= b_0 + b_1 x + b_2 x^2 = B(x), \\
&y^2 - \alpha z^2 + f(x) = c_0 + c_1 x + c_2 x^2 = C(x).
\end{aligned}
\]
Substitute $y = A(x)$ and $z = B(x)$ into $y^2 - \alpha z^2 + f(x) = C(x)$, we have
\[
f(x) + A(x)^2 - \alpha B(x)^2 - C(x) = 0.
\]
Define %the \emph{section polynomial}
\[
P_{\Pi,f}(t) = f(t) + G_\Pi(t), \quad \text{where } G_\Pi(t) = A(t)^2 - \alpha B(t)^2 - C(t).
\]
Then $s_f(x,y,z) \in \Pi$ if and only if $x$ is a root of $P_{\Pi,f}$ over $\mathbb{F}_q$. Moreover, each root $x$ uniquely determines $y=A(x)$ and $z=B(x)$, and then uniquely determines a point in $S_f \cap \Pi$. This gives a bijection between $S_f \cap \Pi$ and the set of distinct roots of $P_{\Pi,f}$ over $\mathbb{F}_q$. In particular, $|S_f \cap \Pi|$ equals the number of distinct roots of $P_{\Pi,f}$ over $\mathbb{F}_q$.

Since $\deg A(t) \leq 2$ and $\deg B(t) \leq 2$, we have $\deg(A(t)^2) \leq 4$ and $\deg(B(t)^2) \leq 4$. So $\deg G_\Pi \leq 4$. By assumption $b \geq 5 > 4$,  the leading term of $P_{\Pi,f}$ is identical to that of $f$. Thus $P_{\Pi,f}$ is always a monic polynomial of degree $b$.



The following claim shows that the subfamily of planes with $b_2 \neq 0$ is closed under the orthogonal polarity.

\begin{claim}\label{clm:polar-closure}
If $b_2 \neq 0$, then $\Pi(A,B,C)^\perp = \Pi(A',B',C')$ for some  polynomials $A',B',C'$ of degree at most $2$.
\end{claim}

\begin{proof}[Proof of Claim \ref{clm:polar-closure}]
Consider the following three points in $\mathbb{F}_q^6$, 
\[
\begin{aligned}
r_A &= (1, 0, -a_0, 0, -a_2, -a_1), \\
r_B &= (0, 0, -b_0, 1, -b_2, -b_1), \\
r_C &= (0, 1, -c_0, 0, -c_2, -c_1).
\end{aligned}
\]
We first verify that each point is orthogonal to every point in $\Pi(A,B,C)$. Take arbitrary $p = (w_p,x_p,y_p,z_p,u_p,v_p) \in \Pi(A,B,C)$. Then
\[
y_p = a_0 w_p + a_1 x_p + a_2 u_p.
\]
Thus, we get
\begin{align*}
\beta(r_A,p)
&= w_{r_A}y_p + y_{r_A}w_p + x_{r_A}v_p + v_{r_A}x_p + z_{r_A}z_p + u_{r_A}u_p \\
&= 1\cdot y_p + (-a_0)\cdot w_p + 0\cdot v_p + (-a_1)\cdot x_p + 0\cdot z_p + (-a_2)\cdot u_p \\
&= y_p - a_0 w_p - a_1 x_p - a_2 u_p\\
&=0
\end{align*}
%Combined with the plane relation $y_p = a_0 w_p + a_1 x_p + a_2 u_p$, we get
%\[
%\beta(r_A,P) = 0.
%\]
Since $p$ is arbitrary in $\Pi(A,B,C)$, we conclude that $r_A \perp \Pi(A,B,C)$. 
The orthogonality of $r_B$ and $r_C$ follows by identical reasoning, using the defining equations for $z_P$ and $v_P$ respectively.

The vectors $r_A, r_B, r_C$ are also linearly independent. Indeed, each has a unique coordinate equal to $1$ (the $w$-coordinate for $r_A$, the $z$-coordinate for $r_B$, and the $x$-coordinate for $r_C$), so no nontrivial linear combination can yield the zero vector. Since $\beta$ is nondegenerate, the orthogonal complement of a $3$-dimensional subspace is also $3$-dimensional. Therefore $\{r_A, r_B, r_C\}$ forms a basis for $\Pi(A,B,C)^\perp$.

Any vector in $\Pi(A,B,C)^\perp$ can be written as $\lambda_1 r_A + \lambda_2 r_B + \lambda_3 r_C$ with $\lambda,\lambda_2,\lambda_3 \in \mathbb{F}_q$ and coordinates
\[
\begin{array}{lll}
w = \lambda_1, &\quad x = \lambda_3, &\quad y = -a_0\lambda_1 - b_0\lambda_2 - c_0\lambda_3, \\
z = \lambda_2, &\quad u = -a_2\lambda_1 - b_2\lambda_2 - c_2\lambda_3, &\quad v = -a_1\lambda_1 - b_1\lambda_2- c_1\lambda_3.
\end{array}
\]
We now express $y,z,v$ as linear functions of $w,x,u$. From the coordinate equations, immediately $\lambda_1 = w$ and $\lambda_3 = x$. 
Substitute these into the expression for $u$, we get that 
\begin{align}\label{eq11}
u = -a_2 w - b_2 \lambda_2 - c_2 x.
\end{align}
Since $b_2 \neq 0$, by (\ref{eq11}), we obtain that 
\begin{align}\label{eq12}
\lambda_2 = -\frac{a_2}{b_2} w - \frac{c_2}{b_2} x - \frac{1}{b_2} u.
\end{align}
Substitute $\lambda_1=w$, $\lambda_3=x$, and the equality (\ref{eq12}) into the formulas for $z$, $y$, and $v$, we obtain that
\[
\begin{aligned}
z &= \left(-\frac{a_2}{b_2}\right) w + \left(-\frac{c_2}{b_2}\right) x + \left(-\frac{1}{b_2}\right) u, \\
y&= \left(-a_0 + \frac{a_2 b_0}{b_2}\right) w
+ \left(-c_0 + \frac{b_0 c_2}{b_2}\right) x
+ \frac{b_0}{b_2} u, \\
v&= \left(-a_1 + \frac{a_2 b_1}{b_2}\right) w
+ \left(-c_1 + \frac{b_1 c_2}{b_2}\right) x
+ \frac{b_1}{b_2} u.
\end{aligned}
\]
All three coordinates $y,z,v$ are linear combinations of $w,x,u$ with constant coefficients. Defining $A',B',C'$ to be the  polynomials of degree at most $2$ with these coefficients, we conclude $\Pi(A,B,C)^\perp = \Pi(A',B',C')$.
\end{proof}


We call a plane $\Pi = \Pi(A,B,C)$ with $b_2 \neq 0$ \emph{good} if both $P_{\Pi,f}$ and $P_{\Pi^\perp,f}$ split completely into distinct linear factors over $\mathbb{F}_q$. Since both are monic polynomials of degree $b$, this condition is equivalent to $|S_f \cap \Pi| = b$ and $|S_f \cap \Pi^\perp| = b$.

Fix an arbitrary plane $\Pi$ in the subfamily. By Claim~\ref{clm:polar-closure}, $\Pi^\perp = \Pi(A',B',C')$ for some $A',B',C'$. Define
\[
G_{\Pi^\perp}(t) = A'(t)^2 - \alpha B'(t)^2 - C'(t),
\]
which satisfies $\deg G_{\Pi^\perp} \leq 4$ by the same degree argument. Let
\[
\delta(t) = G_{\Pi^\perp}(t) - G_\Pi(t).
\]
Then $\deg \delta \leq 4 < b$.

The condition that both $f + G_\Pi$ and $f + G_{\Pi^\perp}$ split completely into distinct linear factors is equivalent to the condition that both $F$ and $F+\delta$ split completely, where we set $F = f + G_\Pi$. Indeed, substituting $G_{\Pi^\perp} = G_\Pi + \delta$ gives
\[
f + G_{\Pi^\perp} = f + G_\Pi + \delta = F + \delta.
\]

We now assert that the map $\Phi: \mathcal{M}_b \to \mathcal{M}_b$ given by $\Phi(f) = f + G_\Pi$ is a bijection.
Since $\deg G_\Pi \leq 4 < b$, adding $G_\Pi$ does not change the leading term of $f$, so $F = f+G_\Pi$ remains monic of degree $b$.
If $\Phi(f_1) = \Phi(f_2)$, then $f_1 + G_\Pi = f_2 + G_\Pi$, so $f_1 = f_2$.
For any $F_0 \in \mathcal{M}_b$, set $f_0 = F_0 - G_\Pi$. 
Again $\deg G_\Pi < b$ implies that $f_0$ is monic of degree $b$, so $f_0 \in \mathcal{M}_b$ and $\Phi(f_0) = F_0$. 
Thus, the assertion holds.

By Lemma~\ref{lem:double-shift}, for fixed $\delta$ with $\deg \delta < b$, the number of $F \in \mathcal{M}_b$ for which both $F$ and $F+\delta$ split completely into distinct linear factors is at least $\eta_b q^b$, where $\eta_b > 0$ depends only on $b$. By the bijection $\Phi$, the number of $f \in \mathcal{M}_b$ making $\Pi$ good is also at least $\eta_b q^b$.
Since $|\mathcal{M}_b| = q^b$, if $f$ is chosen uniformly at random from $\mathcal{M}_b$, the probability that $\Pi$ is good is at least
$\frac{\eta_b q^b}{q^b} = \eta_b.$
This lower bound holds uniformly for every plane in the subfamily.

We next present the expected number of good planes.
Let $X_f$ denote the total number of good planes for a given $f$. 
We count the size of the $b_2 \neq 0$ subfamily:
$A(t)$ has $q^3$ choices (three free coefficients);
$B(t)$ has $q^2(q-1)$ choices, where $b_2 \in \mathbb{F}_q^*$ gives $q-1$ options,
and $b_0,b_1$ each have $q$ options;
and $C(t)$ has $q^3$ choices.
So the total number of planes in the subfamily is
\[
q^3 \cdot q^2(q-1) \cdot q^3 = (q-1) q^8.
\]
By linearity of expectation,
\[
\mathbb{E}_f [X_f] \geq \eta_b \cdot (q-1) q^8.
\]
Since the expectation meets this lower bound, there exists some $f \in \mathcal{M}_b$ for which
\[
X_f \geq \eta_b (q-1) q^8 = \Omega_b(q^9).
\]
That is, there exists an $f$ such that the number of good planes is $\Omega_b(q^9)$.



In the rest of proof, we  convert the count of good planes into a lower bound on the number of $K_{a,b}$ copies in $G_f$.
Let $\Pi$ be a good plane. Then $|S_f \cap \Pi| = b$ and $|S_f \cap \Pi^\perp| = b$. By definition of the orthogonal complement, every point in $\Pi$ is orthogonal to every point in $\Pi^\perp$. Therefore, the left vertices of $L$ corresponding to $S_f \cap \Pi$ and the right vertices of $R$ corresponding to $S_f \cap \Pi^\perp$ induce a complete bipartite subgraph $K_{b,b}$ in $G_f$.

Distinct good planes yield distinct $K_{b,b}$ subgraphs. Suppose two good planes $\Pi_1$ and $\Pi_2$ induce the same left vertex set $L = S_f \cap \Pi_1 = S_f \cap \Pi_2$. Since $b \geq 5$, we may choose any three points from $L$. By Lemma~\ref{lem:no-three-collinear}, no three points of $S_f$ are collinear, so three points uniquely determine a projective plane. Hence $\Pi_1 = \Pi_2$. Thus each good plane corresponds to a unique $K_{b,b}$.

Each such $K_{b,b}$ contains $\binom{b}{a}$ distinct copies of $K_{a,b}$, obtained by selecting any $a$ vertices from the left partition $L$ and all $b$ vertices from the right partition $R$. Since $a \geq 4$, any $K_{a,b}$ obtained this way has at least $3$ left vertices, which uniquely determine the underlying good plane. Therefore,  the number of $K_{a,b}$ copies in $G_f$ satisfies
\[
N(K_{a,b}, G_f) \geq \binom{b}{a} \cdot \eta_b (q-1) q^8 = \Omega_{a,b}(q^9).
\]

The proof of Lemma \ref{lem:good-plane-count} is complete.
\end{proof}

\end{document}
